\documentclass[12pt, a4paper]{article}
\usepackage[english]{babel}
\usepackage{titlesec}

% \titleformat{\chapter} % command [display] % shape {\bfseries\huge} % format
%   {Chapter \thechapter} % label {0pt} % sep {} % before-code


\usepackage[T5]{fontenc}
\usepackage{mathptmx}[ptm]
\usepackage{a4wide, amssymb, epsfig, latexsym, array, hhline, fancyhdr}
\usepackage[normalem]{ulem}
%\usepackage{soul}
\usepackage{float}
\usepackage[makeroom]{cancel}
\usepackage{amsmath}
\usepackage{amsthm}
\usepackage{multicol, longtable, amscd}
\usepackage{diagbox}%Make diagonal lines in tables
\usepackage{booktabs}
\usepackage{alltt}
\usepackage[framemethod=tikz]{mdframed}% For highlighting paragraph backgrounds
\usepackage{indentfirst}
\usepackage{caption,subcaption}
\usepackage{svg}
\graphicspath{ {figures/} }
\usepackage{lastpage}
\usepackage[lined,boxed,commentsnumbered]{algorithm2e}
\usepackage{enumerate}
\usepackage{color}
\usepackage{graphicx}% Standard graphics package
\graphicspath{ {figures/} }
\usepackage{array}
\usepackage{tabularx, caption}
\usepackage{tabularray}
\usepackage{multirow}
\usepackage{multicol}
\usepackage{rotating}
\usepackage{graphics}
\usepackage{geometry}
\usepackage{setspace}
\usepackage{epsfig}
\usepackage{tikz}
\usepackage[numbers]{natbib}
\usepackage{subcaption}

\usepackage{dirtree}



\usetikzlibrary{arrows, snakes, backgrounds, calc}
\usepackage[unicode]{hyperref}
\hypersetup{urlcolor=blue,linkcolor=black,citecolor=black,colorlinks=true} 

\usepackage{xcolor}

\definecolor{codegreen}{rgb}{0,0.6,0}
\definecolor{codegray}{rgb}{0.5,0.5,0.5}
\definecolor{codepurple}{rgb}{0.58,0,0.82}
\definecolor{backcolour}{rgb}{0.95,0.95,0.92}

\usepackage{listings}
\lstdefinestyle{mystyle}{ backgroundcolor=\color{backcolour},   
    commentstyle=\color{codegreen}, keywordstyle=\color{magenta},
    numberstyle=\tiny\color{codegray}, stringstyle=\color{codepurple},
    basicstyle=\ttfamily\footnotesize, breakatwhitespace=false,         
    breaklines=true,                 
    captionpos=b,                    
    keepspaces=true,                 
    numbers=left,                    
    numbersep=5pt,                  
    showspaces=false,                
    showstringspaces=false, showtabs=false,                  
    tabsize=2 }
\lstset{style=mystyle}



\def\thesislayout{	% A4: 210 × 297
	\geometry{ a4paper, total={160mm,240mm},  % fix over page
		left=30mm, top=30mm, } }
\thesislayout

%\usepackage{fancyhdr}
\setlength{\headheight}{40pt}
\pagestyle{fancy}
\fancyhead{} % clear all header fields
\fancyhead[L]{
 \begin{tabular}{rl}
    \begin{picture}(25,15)(0,0) \put(0,-8){\includegraphics[width=10mm,
    height=10mm]{Images/hcmut.png}}
    %\put(0,-8){\epsfig{width=10mm,figure=hcmut.eps}}
   \end{picture}&
	%\includegraphics[width=8mm, height=8mm]{hcmut.png} & %
	\begin{tabular}{l}
        \textbf{ Ho Chi Minh City University of Technology } \\
        \textbf{ Faculty of Computer Science and Engineering }
	\end{tabular} 	
 \end{tabular}
} \fancyhead[R]{
	\begin{tabular}{l}
		\tiny \bf \\
		\tiny \bf \end{tabular}  } \fancyfoot{} % clear all footer fields
\fancyfoot[L]{\scriptsize Capstone Project Report (CO4029) - Academic year 2023 - 2024}
\fancyfoot[R]{\scriptsize  Page {\thepage}/\pageref{LastPage}}
\renewcommand{\headrulewidth}{0.3pt}
\renewcommand{\footrulewidth}{0.3pt}

%%%%%%%%%%%%%%%%%%%%%%%%%%%%%%%%%%%%%%%%%%
% Tạo một môi trường mới cho các bảng đặc tả UC, giúp định dạng kiểu
\newenvironment{usecase_table}[1][Test]
{
\begin{table}[H]
% Left-Right Table padding
    \setlength{\tabcolsep}{6pt}
    \renewcommand{\arraystretch}{1.5}
    % \begin{tabularx}{\textwidth}{bs}
    \def\savedcaption{\caption{#1}}%
    \begin{tabular}{|>{\bfseries} p{0.2\linewidth}|p{0.7\linewidth}|}

} {
    \end{tabular}
    \savedcaption
\end{table}
}

\newenvironment{usecase_enum}
{
    \begin{enumerate*}[itemjoin={\newline}]
} {
    \end{enumerate*}
}

\makeatletter
\newenvironment{acknowledgments}{
	\small
	\begin{center}
	  {\bfseries Acknowledgments\vspace{-.5em}\vspace{\z@}}
	\end{center}
	\quotation
  }{}
\makeatother

% Declaration
\makeatletter
\newenvironment{declaration}{
	\small
	\begin{center}
	  {\bfseries Declaration\vspace{-.5em}\vspace{\z@}}
	\end{center}
	\quotation
  }{}
\makeatother

\makeatletter
\newenvironment{abstr}{
        \small
        \begin{center}
	{\bfseries Abstract \vspace{-.5em}\vspace{\z@}}
        \end{center}
        \quotation
    }{}
\makeatother

\makeatletter
\newenvironment{abstrsss}{
        \small
        \begin{center}
	{\bfseries Abstractss \vspace{-.5em}\vspace{\z@}}
        \end{center}
        \quotation
    }{}
\makeatother
%%%%%%%%%%%%%%%%%%%%%%%%%%%%%%%%%%%%%%%%%%%%%%5
\setcounter{secnumdepth}{4}
\setcounter{tocdepth}{3}
\makeatletter
\newcounter {subsubsubsection}[subsubsection]
\renewcommand\thesubsubsubsection{\thesubsubsection .\@alph\c@subsubsubsection}
\newcommand\subsubsubsection{\@startsection{subsubsubsection}{4}{\z@}%
                                     {-3.25ex\@plus -1ex \@minus -.2ex}%
                                     {1.5ex \@plus .2ex}%
                                     {\normalfont\normalsize\bfseries}}
\newcommand*\l@subsubsubsection{\@dottedtocline{3}{10.0em}{4.1em}}
\newcommand*{\subsubsubsectionmark}[1]{}
\makeatother



\everymath{\color{black}}%make in-line maths symbols blue to read/check easily

\sloppy
\captionsetup[figure]{labelfont={small,bf},textfont={small,it},belowskip=-1pt,aboveskip=-9pt}
%space removed between caption, figure, and text
\captionsetup[table]{labelfont={small,bf},textfont={small,it},belowskip=-1pt,aboveskip=7pt}
\setlength{\floatsep}{5pt plus 2pt minus 2pt}
\setlength{\textfloatsep}{5pt plus 2pt minus 2pt}
\setlength{\intextsep}{10pt plus 2pt minus 2pt}


% Declare a variable named Proc using \newcommand

\newcommand{\Uni}{Ho Chi Minh University - Vietnam National Unversity, Ho Chi Minh City}

\renewcommand{\labelitemi}{\normalfont\bfseries\textendash}
\renewcommand{\labelitemii}{\textbullet}
\renewcommand*\baselinestretch{1.5}\selectfont


\thesislayout

\begin{document}

\begin{titlepage}
    \begin{tikzpicture}[remember picture, overlay]
        \draw[line width = 4pt] ($(current page.north west) + (1.0in,-0.5in)$)
        rectangle ($(current page.south east) + (-0.4in,0.5in)$);
        \draw[line width=1.5pt]
        ($ (current page.north west) + (1.05in,-0.55in) $) rectangle ($ (current
        page.south east) + (-0.45in,0.55in) $);
    \end{tikzpicture}

    \begin{center}
        \large \textbf{VIETNAM NATIONAL UNIVERSITY, HO CHI MINH CITY} \\
        \large \textbf{HO CHI MINH UNIVERSITY OF TECHNOLOGY} \\
        \large \textbf{FACULTY OF COMPUTER SCIENCE AND ENGINEERING}
    \end{center}

    \begin{figure}[h!]
        \begin{center}
            \includegraphics[width=5cm]{Images/hcmut.png}
        \end{center}
    \end{figure}

    \begin{center}
        \begin{tabular}{c}
        \multicolumn{1}{l}{\textbf{{\Large Discrete Structure (CO1007)}}}\\
        ~~\\
        \hline
        \\
        \multicolumn{1}{l}{\textbf{{\Large Assignment Report}}}\\
        \\
        \textbf{{\Huge BELLMAN-FORD ALGORITHM}}\\
        \\
        \hline
        \end{tabular}
        \end{center}

    \vspace{1cm}
    \begin{table}[H]
        \begin{tabular}{rrl}
        \hspace{5 cm} & \textbf{Mentor}: & Nguyen An Khuong \\
        & & Mai Xuan Toan \\
        & & Nguyen Tien Thinh \\
        
        & \textbf{Student}: & Trần Hà Tuấn Kiệt -- 2011493 \\
        
        \end{tabular}
        \end{table}
    \vspace{1cm}

    \begin{center}
        {\large Ho Chi Minh City, \the\month/\the\year}
    \end{center}
\end{titlepage}


\newpage

%\thispagestyle{empty} \chapter*{Work Assignment} \begin{table}[H] \centering
% \renewcommand{\arraystretch}{1.5} \begin{tabular}{|c|c|c|l|} \hline
% \textbf{No.}       & \textbf{Full name}                 & \textbf{Student ID}
% & \textbf{Work Assignment}                     \\
%         \hline %%%%%Student 1%%%%%%%%%%% \multirow{4}{*}{1} &
%         \multirow{4}{*}{Trần Hà Tuấn Kiệt} & \multirow{4}{*}{2011493} & -
%         Exploratory Data Analysis                  \\
%                            &                                    &
%                            & - System proposal and architecture           \\
%                            &                                    &
%                            & - Research and Develop with new technologies \\
%                            &                                    &
%                            & - Implement system endpoints                 \\
%         \hline %%%%%Student 2%%%%%%%%%%% \multirow{4}{*}{2} &
%         \multirow{4}{*}{Nguyễn Đức Thuỵ}   & \multirow{4}{*}{2012158} & -
%         Exploratory Data Analysis                  \\
%                            &                                    &
%                            & - System proposal and architecture           \\
%                            &                                    &
%                            & - Research and Develop with new technologies \\
%                            &                                    &
%                            & - Implement data processing flows            \\
%         \hline \end{tabular} \end{table} \newpage

% Revert title name back to English
\renewcommand{\contentsname}{Table of Contents}
\renewcommand{\listfigurename}{List of Figures}
\renewcommand{\listtablename}{List of Tables}
\renewcommand{\bibname}{References}
\renewcommand{\figurename}{Figure}

\tableofcontents
\newpage

%\thispagestyle{empty}
\listoffigures
\listoftables
\newpage

% \renewcommand\thechapter{Chapter \arabic{chapter}}

\section{Introduction}
Graph algorithms are essential in solving problems related to networks, routing,
and resource allocation. The Bellman-Ford algorithm provides a robust method for
finding shortest paths in graphs that may have negative weights. The Traveling
Salesman Problem (TSP), on the other hand, poses a challenge in combinatorial
optimization, seeking the shortest possible route visiting each city and
returning to the origin city.

\section{Bellman-Ford Algorithm}
\subsection{Description}
The Bellman-Ford algorithm computes shortest paths from a single source vertex
to all other vertices in a weighted digraph. It is capable of handling graphs
with negative weight edges, unlike Dijkstra's algorithm.

\subsection{Complexity Analysis}
The time complexity of Bellman-Ford is \(O(VE)\), where \(V\) is the number of
vertices and \(E\) is the number of edges. Its space complexity is \(O(V)\), as
it maintains a distance array.

\subsection{Bellman-Ford Path Calculation}
Path reconstruction is carried out using the predecessor array filled during the
execution of the Bellman-Ford algorithm. By tracking predecessors from the
destination node back to the source, one can determine the complete path
traversed.

\section{Traveling Salesman Problem}
\subsection{Introduction}
The Traveling Salesman Problem (TSP) presents a classic dilemma in combinatorial
optimization and operations research, asking for the shortest possible route
that visits each city exactly once and returns to the origin city. The problem
is classified as NP-hard, which implies that no efficient algorithm exists that
can solve all instances of TSP within polynomial time. This makes it an ideal
benchmark for many optimization techniques, both exact and heuristic, in fields
ranging from logistics to DNA sequencing, where the goal is often to minimize
travel distance or cost.

\subsection{Theoretical Significance}
Understanding TSP helps illuminate the complexity inherent in many real-world
optimization problems. It also serves as a gateway to discussing other
computational problems in NP, such as the Hamiltonian cycle problem, and
provides insight into the theory of computational complexity itself. Its study
has helped develop foundational concepts in algorithm design and complexity
theory, including the introduction of approximation algorithms for NP-hard
problems.

\subsection{Genetic Algorithm for TSP}
Genetic Algorithms (GAs) are inspired by Darwinian principles of natural
selection and genetics, and they have proven to be particularly adept at finding
good solutions to TSP with relatively low computational costs compared to
exhaustive search methods\cite{8862255}. They simulate the process of natural
selection where the fittest individuals are selected for reproduction in order
to produce offspring of the next generation.

\subsection{Implementation}
GAs for TSP start with a population of possible solutions (tours) represented as
sequences of city visits. These solutions undergo processes analogous to natural
genetic variation\cite{article}: 
\begin{itemize}
    \item \textbf{Selection}: Tours are selected for reproduction based on their
    fitness, which is typically inversely proportional to their travel cost.
    \item \textbf{Crossover}: Selected tours are recombined to produce new
    offspring tours, hoping to combine advantageous traits from each parent.
    \item \textbf{Mutation}: To maintain genetic diversity within the population
    and to explore a wider solution space, tours are randomly altered. 
\end{itemize}
Iterative application of these processes results in the evolution of tours that
tend to be more fit over time.


\subsection{Pseudocode}
Provide detailed pseudocode for a Genetic Algorithm applied to TSP. The
following pseudocode outlines the process of evolving a population of tours to
find an approximate solution to the TSP:

\begin{verbatim}
    function GeneticAlgorithm()
        Initialize population with random tours
        Evaluate fitness of each tour
        Repeat
            Selection: Choose the fitter tours for reproduction
            Crossover: Combine parts of two tours to create new tour
            Mutation: Randomly alter a tour's path
            Evaluate new fitness: Assess the quality of new tours
        Until termination condition met 
        (e.g., number of generations, time limit)
        return best tour found
    end function
\end{verbatim}
    

Each step of the Genetic Algorithm is detailed below:

\subsubsection*{Selection}
The selection process preferentially chooses individuals from the current
population based on their fitness scores.
\begin{verbatim}
Function Selection(Population, FitnessScores)
    Selected = []
    For each tour in Population
        Probability = FitnessScore of tour / Sum of all FitnessScores
        If random number < Probability
            Add tour to Selected
        End If
    End For
    Return Selected
End Function
\end{verbatim}

\subsubsection*{Crossover}
Crossover is a genetic operator used to combine the genetic information of two
parents to generate new offspring.
\begin{verbatim}
Function Crossover(Parent1, Parent2)
    Child = New Tour()
    SplitPoint = Random number between 1 and length of Tour - 1
    Copy first SplitPoint genes from Parent1 to Child
    Fill remaining genes in Child with genes from Parent2 
not already in Child maintaining order
    Return Child
End Function
\end{verbatim}

\subsubsection*{Mutation}
Mutation introduces small random changes in the offspring, providing genetic
diversity.
\begin{verbatim}
Function Mutation(Tour, MutationRate)
    For each gene in Tour
        If random number < MutationRate
            Swap this gene with another random gene in the Tour
        End If
    End For
    Return Tour
End Function
\end{verbatim}

\subsubsection*{Evaluate New Fitness}
After mutation and crossover, the fitness of new solutions must be evaluated.
\begin{verbatim}
Function EvaluateFitness(Tour)
    TotalDistance = 0
    For i from 1 to length of Tour - 1
        TotalDistance += Distance between Tour[i] and Tour[i+1]
    End For
    TotalDistance += Distance between Tour[end] and Tour[begin]
    Fitness = 1 / TotalDistance
    Return Fitness
End Function
\end{verbatim}


\subsection{Results Discussion}
The performance of Genetic Algorithms in solving the TSP can vary significantly
with changes in parameters such as population size, mutation rate, and selection
method. While GAs do not guarantee an optimal solution, they are highly
effective at producing very good solutions within reasonable time limits,
particularly for very large instances where other methods become impractical.
Their effectiveness also hinges on the proper tuning of parameters and
adaptation of the algorithm to specific characteristics of the problem instance.


\section{Evaluation}

\subsection{How to Compile and Run the Program}
To compile and run the program, follow the steps outlined below. These
instructions assume you have a Makefile configured with the necessary commands
to compile the program.

\subsubsection*{Compiling the Program}
To compile the program, open a terminal in the directory containing your source
files and Makefile, and execute the following command:
\begin{verbatim}
make -s
\end{verbatim}
This command calls `make` with the `-s` (silent) option to suppress the display
of commands being executed, making the output cleaner.

\subsubsection*{Running the Program}
Once the compilation is successful, you can run the compiled program by entering
the following command in the terminal:
\begin{verbatim}
./main
\end{verbatim}
This command executes the `main` executable that was generated by the
compilation process. Ensure that you have the appropriate permissions to execute
the file, and that the current working directory is where the `main` executable
is located.


\subsection{Test Results}

\[
\begin{array}{cccccccccccc}
0 & 17 & 40 & 100 & 32 & 11 & 29 & 77 & 73 & 53 & 52 & 72 \\
48 & 0 & 23 & 42 & 81 & 5 & 12 & 37 & 31 & 55 & 18 & 66 \\
48 & 97 & 0 & 74 & 21 & 43 & 8 & 37 & 48 & 13 & 55 & 88 \\
79 & 53 & 66 & 0 & 8 & 97 & 38 & 61 & 48 & 76 & 75 & 71 \\
18 & 56 & 42 & 97 & 0 & 59 & 40 & 51 & 77 & 5 & 65 & 39 \\
45 & 86 & 82 & 20 & 23 & 0 & 96 & 100 & 44 & 84 & 45 & 97 \\
15 & 53 & 93 & 53 & 80 & 7 & 0 & 29 & 21 & 78 & 13 & 43 \\
87 & 9 & 1 & 26 & 26 & 44 & 98 & 0 & 92 & 49 & 9 & 77 \\
30 & 29 & 66 & 93 & 28 & 9 & 76 & 73 & 0 & 73 & 91 & 92 \\
31 & 9 & 38 & 38 & 38 & 25 & 83 & 17 & 68 & 0 & 69 & 93 \\
35 & 61 & 19 & 45 & 24 & 10 & 94 & 33 & 53 & 90 & 0 & 62 \\
19 & 82 & 57 & 95 & 57 & 96 & 67 & 13 & 54 & 65 & 90 & 0 \\
\end{array}
\]

\subsubsection*{TSP}

\textbf{Parameters}: Population size = 100, Mutation Rate = 0.05, Generations =
2000

\textbf{Results}:
\begin{verbatim}
    Running TSP Solver:
    TSP Path: IFDEJABKLHCGI
    TSP Path Cost: 213
    Accuracy: 95.5882%
\end{verbatim}

\textbf{Parameters}: Population size = 200, Mutation Rate = 0.05, Generations =
2000

\textbf{Results}:
\begin{verbatim}
    Running TSP Solver:
    TSP Path: GIAFDEJBKLHCG
    TSP Path Cost: 206
    Accuracy: 99.0196%
\end{verbatim}

\textbf{Parameters}: Population size = 200, Mutation Rate = 0.05, Generations =
10000

\textbf{Results}:
\begin{verbatim}
    Running TSP Solver:
    TSP Path: HKCJBGIAFDELH
    TSP Path Cost: 204
    Accuracy: 100%
\end{verbatim}

\textbf{Parameters}: Population size = 200, Mutation Rate = 0.07, Generations =
2000

\textbf{Results}:
\begin{verbatim}
    Running TSP Solver:
    TSP Path: GIAFDELHKCJBG
    TSP Path Cost: 204
    Accuracy: 100%
\end{verbatim}

\subsection{Limitations and Potential Improvements}
Despite their advantages, GAs can suffer from premature convergence to
suboptimal solutions and can be sensitive to the initial population and
parameter settings. Ongoing research aims to improve the robustness and
efficiency of GAs by integrating advanced techniques such as hybrid algorithms
that combine GAs with local search methods, adaptive parameter tuning, and
machine learning strategies to predict and adjust evolutionary dynamics
dynamically.

\section{Conclusion}
The Bellman-Ford algorithm's ability to handle negative weights makes it a
powerful tool for graph analysis. Meanwhile, GAs offer a robust heuristic method
for tackling NP-hard problems like the TSP, especially in large datasets where
exact methods are computationally impractical.


\bibliographystyle{plainnat}
\bibliography{refs}

\end{document}